\section{Pytania badawcze} \label{sec:rq}

W pracy postawiono cztery pytania badawcze, którym podporządkowano projekt eksperymentów oraz strukturę analizy wyników w rozdziale~\ref{chap:results}.

\textbf{RQ1: Czy rozszerzenie funkcji oceny o siedem nowych cech poprawia skuteczność agenta?}

Przejście z 21 do 28 parametrów wprowadza cechy specyficzne dla konkretnych talii i mechanik (trwałość broni, synergia piratów, postęp Jade Golem), których agent bazowy nie uwzględniał. Pytanie dotyczy tego, czy te dodatkowe informacje są na tyle ważne strategicznie, aby przełożyć się na wyższy wskaźnik wygranych.

\textbf{RQ2: Czy podział oceny na trzy fazy gry i rozszerzenie reprezentacji do 63 parametrów poprawia skuteczność agenta?}

Agent trójfazowy zakłada, że optymalne zachowanie różni się strategicznie pomiędzy wczesną, środkową i późną fazą rozgrywki. Pytanie bada, czy niezależna optymalizacja wag dla każdej fazy pozwala agentowi lepiej dopasować strategię do kontekstu tury, czy też złożoność przestrzeni poszukiwań $[0,1]^{63}$ stanowi zbyt duże wyzwanie dla dostępnego budżetu obliczeniowego.

\textbf{RQ3: Czy strategia przeszukiwania głębokościowego (\textit{lookahead}) poprawia skuteczność w porównaniu do wariantu bazowego?}

Wariant głębokościowy zachowuje identyczny zestaw 21 wag, ale zmienia horyzont decyzyjny z 1 na 2 akcje. Pytanie sprawdza, czy ta zmiana strategii podejmowania decyzji --- bez żadnego zwiększenia liczby optymalizowanych parametrów --- jest opłacalna i czy uzyskuje istotną przewagę nad agentem bazowym.

\textbf{RQ4: Który optymalizator --- koewolucja, SHADE koewolucyjny czy SHADE z Hall of Fame --- wytrenuje najskuteczniejszych agentów przy równoważnym budżecie obliczeniowym?}

Algorytm SHADE wprowadza adaptacyjną regulację parametrów $F$ i $CR$ na podstawie historii sukcesów poprzednich pokoleń. Pytanie skupia się na tym, czy to podejście przekłada się na wyższy wskaźnik
wygranych agenta końcowego w porównaniu z wariantem koewolucyjnym. Oceny dokonuje się poprzez zestawienie wyników mini-turniejów Fazy~I, w których dla każdego wariantu reprezentacji bezpośrednio rywalizują ze sobą najlepsi kandydaci uzyskani przez każdy z trzech optymalizatorów.

Odpowiedzi na pytania RQ1--RQ3 formułowane będą na podstawie wyników turnieju głównego (sekcja~\ref{sec:phase2}). Odpowiedź na RQ4 wynikać będzie natomiast z analizy Fazy~I (sekcja~\ref{sec:phase1}). Dla każdego wariantu reprezentacji porównuje się, który z trzech optymalizatorów --- koewolucja, SHADE koewolucyjny oraz SHADE z Hall of Fame --- okazał się najsilniejszym reprezentantem.
